\section{Evaluation}\label{sec:hub:eval}

\paragraph{Free-running decoding.} The hypothesis is tested under free-running decoding.
Teacher forcing would hand the decoder the true previous tokens, and that prefix carries information about the next token.
Under free-running decoding, each decoder instead conditions on the soft prompt and on its own generated prefix.
Greedy generation starts at the start token and stops at the end-of-sequence token or at the generation cap.
The cap is the length of the soft prompt plus \vf{8} tokens on NQ and plus \vf{32} tokens on the legal datasets. % src: configs/v6_3.yaml (cap_margin 8); legal hub configs (32); NQ re-evaluated at 32 moves no direction by more than 0.001 (commit e1ec486)

\paragraph{Metrics.} Two metrics compare the reconstruction with the source prompt.
Word-F1 is the harmonic mean of word precision and recall over word multisets, averaged over prompts.
Exact match is the fraction of reconstructions equal to their source prompt after Unicode and whitespace normalization.
Word-F1 ignores word order, so exact match is reported beside it as the order-sensitive metric.

\paragraph{Mismatched-prompt control.} A fluent decoder may score above zero without reading its soft prompt, because prompts of one dataset share common words.
The mismatched-prompt control measures that floor: the decoder receives the hub outputs of another prompt, and its output is scored against the original prompt.
Prompt $i$ receives the hub outputs of prompt $i+1$, and the last prompt those of the first. % src: src/v6_3/inversion.py (shuffled_indices: (i + 1) mod n); src/v8/projectors.py (wrong_partner: same rule)
A score above its control shows that the reconstruction depends on the prompt's own hub outputs.

\paragraph{NQ.} The NQ dataset contains query strings from Natural Questions, split into \vf{40{,}000} training, \vf{5{,}000} validation, and \vf{5{,}000} test prompts.
Near-duplicate questions never cross a split. % src: configs/v3.yaml (split_schema_version v2-stage2-textcluster-v1); src/v3/splits.py
The validation prompts average \vf{9} words and range from \vf{6} to \vf{19}. % src: configs/v3.yaml (n_train, n_val, n_test); word counts from the references in results/v6_3/hpc-final/inversion-records.jsonl (mean 9.08, median 9)
All scores in this chapter are computed on the validation split, which also selects the checkpoint, so no score is held out; the test split has not been evaluated.

\paragraph{Summary score.} One score summarizes a hub: mean cross-model word-F1, the free-running word-F1 averaged over the six cross-model directions on the validation split.
The hypothesis is tested by that mean against the mismatched-prompt control on the same prompts.

